\subsection{环的直分解}

设 A 为一个环，且 \((b_{i})_{i\in I}\) 为 A 的一族双边理想。我们将称同态

\[
x \mapsto (\phi_ {i} (x)) _ {i \in I},
\]

其中 \(\phi_{i}\) 是 A 到 \(\mathbf{A} / \mathfrak{b}_i\) 上的典范同态，即 A 到 \(\prod_{i\in I}(\mathbf{A} / \mathfrak{b}_i)\) 中的典范同态。

命题 9. 设 A 是一个环，且 \((\mathfrak{b}_1, \ldots, \mathfrak{b}_n)\) 是 A 的双边理想，满足 \(\mathfrak{b}_i + \mathfrak{b}_j = \mathbf{A}\) （对 \(i \neq j\) ）。A 到 \(\prod_{i=1}^{n} (\mathbf{A} / \mathfrak{b}_i)\) 中的典范同态是满射的，其核为 \(\bigcap_{i=1}^{n} \mathfrak{b}_i = \sum_{\sigma \in \mathfrak{E}_n} \mathfrak{b}_{\sigma(1)} \mathfrak{b}_{\sigma(2)} \ldots \mathfrak{b}_{\sigma(n)}\) 。

显然该核为 \(\bigcap_{i=1}^{n} \mathfrak{b}_i\) 。为证明满射性，需要证明：对每一族 \((x_i)_{1 \leqslant i \leqslant n}\) （由 A 的元素组成），存在 \(x \in \mathbf{A}\) ，使得 \(x \equiv x_i (\mathfrak{b}_i)\) 对所有 \(1 \leqslant i \leqslant n\) 成立。我们对 I 的基数 \(n\) 作归纳来证明这一断言， \(n \leqslant 1\) 的情形是平凡的。根据归纳假设，存在 \(y \in \mathbf{A}\) ，使得 \(y \equiv x_i (\mathfrak{b}_i)\) （对 \(1 \leqslant i \leqslant n-1\) ）。我们寻找一个 \(x\) ，其

形式 \(y + z\) （其中 \(z \in \mathbf{A}\) ）。必然地， \(z \equiv 0(\mathfrak{b}_i)\) （对 \(i < n\) ），即 \(z \in \mathfrak{b} = \bigcap_{i=1}^{\infty} \mathfrak{b}_i\) ，而另一方面 \(z \equiv x_n - y(\mathfrak{b}_n)\) 。现在由第 9 号命题 6 有 \(\mathfrak{b}_n + \mathfrak{b} = A\) ，由此可知 \(z\) 存在。最后，核的第二种表达式由第 9 号命题 7 得出。

定义 7. 设 A 为一个环。由 A 的双边理想组成的有限族 \((\mathfrak{b}_i)_{i\in I}\) ，如果 A 到 \(\prod_{i\in I}(\mathbf{A} / \mathfrak{b}_i)\) 中的典范同态是一个同构，就称为 A 的一个直分解。

命题 10. 设 A 为一个环，A' 为其中心，且 \((\mathfrak{b}_{i})_{i\in I}\) 为 A 的一个有限双边理想族。以下条件是等价的：

\begin{enumerate}
\def\labelenumi{(\alph{enumi})}
\item
  族 \((\mathfrak{b}_i)_{i\in I}\) 是 A 的一个直分解；
\item
  存在一个族 \((e_i)_{i\in I}\) （由 \(\mathbf{A}'\) 的幂等元组成），使得 \(e_i e_j = 0\) （对 \(i\neq j\) ）， \(1 = \sum_{i\in I}e_i\) 且 \(\mathfrak{b}_i = \mathbf{A}(1 - e_i)\) （对 \(i\in I\) ）；
\item
  \(\mathfrak{b}_i + \mathfrak{b}_j = \mathbf{A}\) （对 \(i \neq j\) ）且 \(\bigcap_{i \in I} \mathfrak{b}_i = \{0\}\) ；
\item
  \(\mathfrak{b}_i + \mathfrak{b}_j = \mathbf{A}\) （对 \(i \neq j\) ）且 \(\prod_{i \in I} \mathfrak{b}_i = \{0\}\) （对 \(\mathbf{I}\) 上的每一个全序）；
\item
  存在一个直分解 \((\mathfrak{b}_i')_{i\in \mathbf{I}}\) （\(\mathbf{A}'\) 的），使得 \(\mathfrak{b}_i = \mathbf{Ab}_i'\) （对 \(i\in \mathbf{I}\) ）。
\end{enumerate}

(a) \(\Rightarrow\) (b)。若条件 (a) 成立，则 A 可与环 \(\prod_{i\in I}(\mathbf{A} / b_i)\)

以及 \(b_{i}\) 与 \(pr_{i}\) 的核等同。具有 (b) 中性质的 \(e_{i}\) 的存在性则由第 10 号得出。

\begin{enumerate}
\def\labelenumi{(\alph{enumi})}
\setcounter{enumi}{1}
\tightlist
\item
  \(\Rightarrow\) (d). 假设 \(e_i\) 存在且具有性质 (b)。对 \(i \neq j\) ， \(1 - e_i \in \mathfrak{b}_i\) ， \(e_i = e_i(1 - e_j) \in \mathfrak{b}_j\) ，因此 \(1 \in \mathfrak{b}_i + \mathfrak{b}_j\) 且 \(\mathbf{A} = \mathfrak{b}_i + \mathfrak{b}_j\) 。另一方面，若给 I 赋予一个全序且 \((x_i)_{i \in I}\) 是 A 中元素的一个族，则由于 \(e_i\) 是中心的，
\end{enumerate}

\[
\prod_ {i \in \mathrm{I}} x _ {i} (1 - e _ {i}) = \Bigl (\prod_ {i \in \mathrm{I}} x _ {i} \Bigr) \Bigl (\prod_ {i \in \mathrm{I}} (1 - e _ {i}) \Bigr) = \Bigl (\prod_ {i \in \mathrm{I}} x _ {i} \Bigr) \Bigl (1 - \prod_ {i \in \mathrm{I}} e _ {i} \Bigr) = 0
\]

因此 \(\prod_{i\in I}b_i = \{0\}\) 。

\begin{enumerate}
\def\labelenumi{(\alph{enumi})}
\setcounter{enumi}{3}
\item
  \(\Rightarrow\) (c)。这由第 9 号命题 7 得出。
\end{enumerate}

(c) \(\Rightarrow\) (a)。这由命题 9 得出。

因此条件 (a)、(b)、(c) 和 (d) 是等价的。假设它们成立。由 \((\mathbf{b}) \Rightarrow (\mathbf{a})\) ，族 \(\mathfrak{b}_i' = \mathbf{A}'(1 - e_i)\) 是 \(\mathbf{A}'\) 的一个直分解。则 \(\mathfrak{b}_i = \mathbf{A}(1 - e_i) = \mathbf{Ab}_i'\) 对所有 \(i \in \mathbf{I}\) 成立。因此条件 (e) 成立。

最后，假设条件 (e) 成立。由 (a) \(\Rightarrow\) (b)，存在一个族 \((e_i)_{i\in I}\) （由 \(A'\) 的幂等元组成），使得 \(e_i e_j = 0\) （对 \(i \neq j\) ）， \(1 = \sum_{i\in I} e_i\) 且 \(\mathfrak{b}_i' = A'(1 - e_i)\) （对 \(i \in I\) ）。则 \(\mathfrak{b}_i = Abi' = A(1 - e_i)\) （对 \(i \in I\) ），因此条件 (b) 成立。

注。设 A 是一个环。设 \((a_{i})_{i\in I}\) 是 A 的加法群 \(A^{+}\) 的一个有限子群族，使得 \(A^{+}\) 是诸 \(a_{i}\) 的直和。假设 \(a_{i}a_{i}\subset ai\) （对 \(i\in I\) ）且 \(a_{i}a_{j} = \{0\}\) （对 \(i\neq j\) ）。则 \(a_{i}\) 对所有 \(i\in I\) 是 A 的双边理想。以由 A 上的加法和乘法诱导的运算， \(a_{i}\) 是一个环，其单位元是 \(1\in A\) 在 \(a_{i}\) 中的分量。若 \(b_{i} = \sum_{j\neq i}a_{j}\) ，显然诸 \(b_{i}\) 满足命题 10 的条件 (c)，因此 \((b_{i})_{i\in I}\) 是 A 的一个直分解，称其由 \((a_{i})_{i\in I}\) 所定义。

\BourbakiRunInHeading{例：Z的理想与商环}

\(\mathbf{Z}\) 的一个理想是 \(\mathbf{Z}\) 的加法子群，因此具有形式 \(n\) . \(\mathbf{Z}\) （其中 \(n \geqslant 0\) ）；反之，对每个整数 \(n \geqslant 0\) ，集合 \(n\) . \(\mathbf{Z}\) 是一个理想，即主理想 \((n)\) 。因此 \(\mathbf{Z}\) 的每个理想都是主理想，且可唯一地表示为 \(n\mathbf{Z}\) （其中 \(n \geqslant 0\) ）的形式。理想 (1) 等于 \(\mathbf{Z}\) ，理想 (0) 由 0 组成，因而不同于 \(\mathbf{Z}\) 和 \(\{0\}\) 的理想具有形式 \(n\mathbf{Z}\) （其中 \(n > 1\) ）。若 \(m \geqslant 1\) 且 \(n \geqslant 1\) ，则 \(m\mathbf{Z} \supset n\mathbf{Z}\) 当且仅当 \(n \in m\) . \(\mathbf{Z}\) ，即 \(m\) 整除 \(n\) 。因此，理想 \(n\mathbf{Z}\) 为极大的必要且充分条件是：不存在整数 \(m > 1\) （异于 \(n\) 且整除 \(n\) ）；换言之， \(\mathbf{Z}\) 的极大理想正是形如 \(p\mathbf{Z}\) 的理想，其中 \(p\) 是一个素数（§ 4，第 10 号，定义 16）。

设 \(m\) 和 \(n\) 是两个整数 \(\geqslant 1\) 。理想 \(m\mathbf{Z} + n\mathbf{Z}\) 是主理想，因此存在一个整数 \(d \geqslant 1\) ，以 \(d\mathbf{Z} = m\mathbf{Z} + n\mathbf{Z}\) 为特征；对每个整数 \(r \geqslant 1\) ，关系“r 整除 d”等价于 \(rZ \supset dZ\) ，从而等价于“ \(rZ \supset mZ\) 且 \(rZ \supset nZ\) ”，即“r 整除 m 和 n”。由此看出，m 和 n 的公因数正是 d 的因数，且 d 是 m 和 n 公有的因数中 \(\geqslant 1\) 者里最大的；d 称为 m 和 n 的最大公因数（缩写为 g.c.d.）。由于 \(dZ = mZ + nZ\) ，存在两个整数 x 和 y，使得 \(d = mx + ny\) 。若 m 和 n 的最大公因数等于 1，则称 m 和 n 互素。这相当于假设存在整数 x 和 y，使得 \(mx + ny = 1\) 。

理想 mZ 与 nZ 的交集非零，因为它包含 mn，因此具有 rZ 的形式，其中 \(r \geqslant 1\) 。如上论证可以看出，r 的倍数正是 m 与 n 的公倍数，且 r 是 m 与 n 的公倍数中 \(\geqslant 1\) 的整数里最小的一个；它称为 m 与 n 的最小公倍数（缩写为 l.c.m.）。

理想 \(m\mathbf{Z}\) 和 \(n\mathbf{Z}\) 的乘积是由形如 \(\sum_{i=1}^{r} mx_i ny_i = mn\left(\sum_{i=1}^{r} x_i y_i\right)\) （对 \(x_1, \ldots, y_r \in \mathbf{Z}\) ）的和组成的集合，因此等于 \(mn\mathbf{Z}\) 。

对每个整数 \(n \geqslant 1\) ，商环 Z/nZ 称为模 n 的整数环；它有 n 个元素，即整数 0, 1, 2, \(\ldots\) , n - 1 模 n 的类。当 n = 1 时，我们得到零环。

命题 11. 设 \(n_1, \ldots, n_d\) 是两两互素的整数 \(\geqslant 1\) ，且 \(n = n_1 \ldots n_r\) 。 \(\mathbf{Z}\) 到乘积环 \(\prod_{i=1}^{r} \mathbf{Z} / n_i \mathbf{Z}\) 中的典范同态是满射的，其核为 \(n\mathbf{Z}\) ，并定义了 \(\mathbf{Z} / n\mathbf{Z}\) 到 \(\prod_{i=1}^{r} \mathbf{Z} / n_i \mathbf{Z}\) 上的一个环同构。

设 \(a_{i} = n_{i}\mathbf{Z}\) （对 \(i = 1, \ldots, r\) ）。根据假设， \(a_{i} + a_{j} = \mathbf{Z}\) （对 \(i \neq j\) ）。于是该命题由命题 9 得出。

上述结果，连同关于素因数分解的那些结果，将在第七章，§ 1 中加以推广，该节专门研究主理想整环；并在《交换代数》第七章，§ 3 中加以推广，该节专门研究阶乘整环。
% semantic-fragments: legacy-input-seam
\relax{}
